\section{Introduction}

Graph classification is a fundamental task in graph representation learning, with widespread applications in bioinformatics (molecular property prediction), cheminformatics (drug discovery), social network analysis (community detection), and computer vision (scene graph understanding). Unlike node-level prediction, graph classification requires aggregating information from entire graph structures to predict labels at the graph level. This task presents unique challenges: graphs exhibit structural heterogeneity (varying sizes, topologies, and connectivity patterns), contain noisy or incomplete information, suffer from long-range dependencies between distant nodes, and face training difficulties in deep architectures due to over-smoothing and gradient degradation. These challenges necessitate models that can effectively capture multi-scale structural features while maintaining trainability at depth.

Graph Neural Networks (GNNs) have emerged as a powerful paradigm for graph representation learning by recursively aggregating neighborhood information. Various convolution operators have been proposed, each with distinct inductive biases: Graph Convolutional Networks (GCN) \cite{kipf2016semi} perform low-pass filtering through normalized adjacency matrices, Graph Attention Networks (GAT) \cite{velivckovic2017graph} adaptively weight neighbor contributions via attention mechanisms, GraphSAGE \cite{hamilton2017inductive} samples and aggregates features from local neighborhoods, and Graph Transformers \cite{dwivedi2020generalization} apply self-attention to capture long-range dependencies. For graph-level prediction, these base operators are typically stacked with pooling layers (e.g., DiffPool \cite{ying2018hierarchical}, SAGPool \cite{lee2019self}) to coarsen graphs hierarchically, followed by readout functions that produce fixed-size graph representations. Despite their success, these approaches face a \textit{unified bottleneck}: as depth increases, repeated neighborhood aggregation causes node representations to converge toward indistinguishable values (over-smoothing \cite{li2018deeper, chen2023revisiting}), while the low-pass nature of most operators filters out high-frequency structural information essential for discriminative power \cite{xu2018powerful}.

To address these limitations, prior work has explored several directions. \textbf{Deep training techniques} such as residual connections \cite{he2016deep, bresson2017residual}, normalization (PairNorm \cite{zhao2020pair}, GraphNorm \cite{cai2021graphnorm}), and edge dropout \cite{rong2019dropedge} alleviate over-smoothing but typically preserve sequential layer-wise architectures. \textbf{Multi-scale architectures} like Jumping Knowledge (JK) networks \cite{xu2018representation} adaptively combine representations from different depths, while DenseGNN \cite{du2024densegnn} connects all layers densely. However, these methods primarily operate within a single branch and do not explicitly model cross-branch information exchange. \textbf{Operator fusion} approaches attempt to combine multiple aggregation schemes (e.g., MixHop \cite{abu2019mixhop} mixes neighborhoods of different sizes), but they often treat operators as independent modules rather than jointly optimized components. \textbf{Pooling-based methods} \cite{ying2018hierarchical, lee2019self} create hierarchical representations through node clustering, yet they risk losing information early when discarding nodes. Despite these advances, a critical gap remains: existing frameworks lack a \textit{unified mechanism} to flexibly integrate diverse graph operators while maintaining robust information flow across multiple branches and depth scales through structured cross-branch residual connections.

In this work, we propose Enhanced Cross Residual Graph Neural Networks (ECR-GNN), a novel framework that unifies multiple graph operators through cross-branch residual connections for improved graph classification. Unlike traditional GNNs that stack layers sequentially or simply add skip connections, ECR-GNN constructs multiple parallel branches—each employing different base operators (GCN, GAT, or Transformer)—and enables them to exchange information through carefully designed cross-residual pathways. Specifically, our approach introduces two key mechanisms: (1) \textit{node-level cross-residual connections}, where intermediate representations from one branch are injected into downstream layers of another branch at each depth step, and (2) \textit{graph-level residual propagation}, where graph-level embeddings from previous branch outputs are added as auxiliary input to subsequent branches. As observed in our implementation on TUDataset benchmarks (MUTAG, DD, MSRC\_9, AIDS), this cross-branch communication allows each operator to leverage complementary inductive biases—for example, GCN's smoothing properties help GAT stabilize attention weights, while GAT's selectivity refines GCN's low-pass features. The framework is implemented modularly using PyTorch Geometric \cite{fey2019fast}, with base operators dynamically instantiated via a unified interface, enabling systematic evaluation of different operator combinations.

The contributions of this work are summarized as follows:

\begin{itemize}
    \item We propose Enhanced Cross Residual Graph Neural Networks (ECR-GNN), a unified framework that integrates multiple graph operators (GCN, GAT, Transformer) through structured cross-branch residual connections, enabling flexible operator combination while preserving trainability at depth.
    \item We design two complementary cross-residual mechanisms: node-level connections that exchange intermediate representations between branches at each layer, and graph-level connections that propagate branch outputs as auxiliary inputs to subsequent branches, facilitating robust information flow across operators and depth scales.
    \item We conduct comprehensive experiments on TUDataset benchmarks, demonstrating that ECR-GNN achieves competitive performance compared to baseline models (BlockGNN, ResBlockGNN, GraphBlockGnn) across different depth configurations (1-5 layers) and embedding dimensions (8, 16, 32, 64), with cross-branch architectures consistently outperforming single-branch variants.
    \item We provide an open-source implementation built on PyTorch Geometric that supports modular operator instantiation and flexible branching configurations, facilitating future research on cross-branch GNN architectures for graph classification tasks.
\end{itemize}
